% Pré-preenchimento: revisar antes da entrega.
\section{Conclusão e Trabalhos Futuros}
\label{sec:conclusao}

\textcolor{orange}{Este trabalho investigou se a avaliação da qualidade de descrições de ferramentas MCP é influenciada pela disponibilização do código da ferramenta ao avaliador. Para isso, 12.171 ferramentas de 206 servidores MCP foram extraídas, junto com o código alcançado por sua árvore de chamadas, e avaliadas por três LLMs segundo a rubrica de seis componentes de Hasan et al. (2026)\nocite{2026arXiv260214878M}, com e sem acesso ao código.}

\textcolor{orange}{Quanto à RQ1, na maior parte da amostra de divergência confirmada (89,2\%) a mudança de nota decorre de um achado concreto no código, sobretudo parâmetros e opções não documentados, modos de falha não divulgados e detalhes da estrutura de saída; isso é consistente com um ganho de precisão, embora não o comprove sem um padrão-ouro humano. Quanto à RQ2, a direção do efeito depende do juiz: o DeepSeek tende a reduzir as notas, enquanto o Gemini e o Qwen3 tendem a elevá-las; \textit{Purpose} é o atributo mais afetado e \textit{Limitations} e \textit{Examples} os menos afetados. Quanto à RQ3, as diferenças são estatisticamente significativas nos seis atributos, com mudança de faixa de quartil entre 31,8\% e 39,1\% das ferramentas em \textit{Purpose}, \textit{Parameter Explanation}, \textit{Length \& Completeness} e \textit{Guidelines}.}

\textcolor{orange}{Conclui-se que o código altera de forma consistente a avaliação das descrições e que essa alteração, na maioria dos casos, tem origem em informações que a descrição omite. A principal contribuição do trabalho é mostrar que avaliar descrições MCP isoladamente deixa de captar uma parcela relevante de suas deficiências, e oferecer uma caracterização empírica, em 16 categorias, do tipo de informação que o código revela.}

\textcolor{orange}{Como trabalhos futuros, propõe-se: (i) construir um padrão-ouro humano sobre uma amostra de ferramentas, permitindo medir diretamente se a avaliação com código é mais precisa; (ii) verificar se as deficiências reveladas pelo código afetam o desempenho real de agentes, por exemplo em \textit{benchmarks} baseados em execução como o MCP-Universe \cite{2025arXiv250814704L}; (iii) usar o código para reescrever automaticamente as descrições e avaliar o ganho de qualidade resultante; (iv) medir a concordância entre avaliadores na análise qualitativa, com um segundo codificador independente e o coeficiente Kappa por categoria; e (v) aprimorar a extração do código, com resolução de chamadas baseada em tipos e profundidade adaptativa.}
